\section{Finite projective modules}
\label{section-finite-projective-modules}

\begin{definition}
\label{definition-locally-free}
Let $R$ be a ring and $M$ an $R$-module.
\begin{enumerate}
\item We say that $M$ is {\it locally free} if we can cover $\Spec(R)$ by
standard opens $D(f_i)$, $i \in I$ such that $M_{f_i}$ is a free
$R_{f_i}$-module for all $i \in I$.
\item We say that $M$ is {\it finite locally free} if we can choose
the covering such that each $M_{f_i}$ is finite free.
\item We say that $M$ is {\it finite locally free of rank $r$}
if we can choose the covering such that each $M_{f_i}$ is isomorphic
to $R_{f_i}^{\oplus r}$.
\end{enumerate}
\end{definition}

\noindent
Note that a finite locally free $R$-module is automatically
finitely presented by Lemma \ref{lemma-cover}. Moreover, if $M$ is a
finite locally free module of rank $r$ over a ring $R$ and if
$R$ is nonzero, then $r$ is uniquely determined by
Lemma \ref{lemma-rank} (because at least one of the
localizations $	R_{f_i}$ is a nonzero ring).

\begin{lemma}
\label{lemma-finite-projective}
Let $R$ be a ring and let $M$ be an $R$-module.
The following are equivalent
\begin{enumerate}
\item $M$ is finitely presented and $R$-flat,
\item $M$ is finite projective,
\item $M$ is a direct summand of a finite free $R$-module,
\item $M$ is finitely presented and
for all $\mathfrak p \in \Spec(R)$ the
localization $M_{\mathfrak p}$ is free,
\item $M$ is finitely presented and
for all maximal ideals $\mathfrak m \subset R$ the
localization $M_{\mathfrak m}$ is free,
\item $M$ is finite and locally free,
\item $M$ is finite locally free, and
\item $M$ is finite, for every prime $\mathfrak p$ the module
$M_{\mathfrak p}$ is free, and the function
$$
\rho_M : \Spec(R) \to \mathbf{Z}, \quad
\mathfrak p
\longmapsto
\dim_{\kappa(\mathfrak p)} M \otimes_R \kappa(\mathfrak p)
$$
is locally constant in the Zariski topology.
\end{enumerate}
\end{lemma}

\begin{proof}
First suppose $M$ is finite projective, i.e., (2) holds.
Take a surjection $R^n \to M$ and let $K$ be the kernel.
Since $M$ is projective,
$0 \to K \to R^n \to M \to 0$ splits.
Hence (2) $\Rightarrow$ (3).
The implication (3) $\Rightarrow$ (2) follows from the fact that
a direct summand of a projective is projective, see
Lemma \ref{lemma-characterize-projective}.

\medskip\noindent
Assume (3), so we can write $K \oplus M \cong R^{\oplus n}$.
So $K$ is a  direct summand of $R^n$ and thus finitely generated.
This shows $M = R^{\oplus n}/K$ is finitely presented.
In other words, (3) $\Rightarrow$ (1).

\medskip\noindent
Assume $M$ is finitely presented and flat, i.e., (1) holds.
We will prove that (7) holds. Pick any prime $\mathfrak p$ and
$x_1, \ldots, x_r \in M$ which map to a basis of
$M \otimes_R \kappa(\mathfrak p)$. By
Nakayama's lemma (in the form of Lemma \ref{lemma-NAK-localization})
these elements generate $M_g$ for some $g \in R$, $g \not \in \mathfrak p$.
The corresponding surjection $\varphi : R_g^{\oplus r} \to M_g$
has the following two properties: (a) $\Ker(\varphi)$ is a finite
$R_g$-module (see Lemma \ref{lemma-extension})
and (b) $\Ker(\varphi) \otimes \kappa(\mathfrak p) = 0$
by flatness of $M_g$ over $R_g$ (see
Lemma \ref{lemma-flat-tor-zero}).
Hence by Nakayama's lemma again there exists a $g' \in R_g$ such that
$\Ker(\varphi)_{g'} = 0$. In other words, $M_{gg'}$ is free.

\medskip\noindent
A finite locally free module is a finite module, see
Lemma \ref{lemma-cover},
hence (7) $\Rightarrow$ (6).
It is clear that (6) $\Rightarrow$ (7) and that (7) $\Rightarrow$ (8).

\medskip\noindent
A finite locally free module is a finitely presented module, see
Lemma \ref{lemma-cover},
hence (7) $\Rightarrow$ (4).
Of course (4) implies (5).
Since we may check flatness locally (see
Lemma \ref{lemma-flat-localization})
we conclude that (5) implies (1).
At this point we have
$$
\xymatrix{
(2) \ar@{<=>}[r] & (3) \ar@{=>}[r] & (1) \ar@{=>}[r] &
(7)  \ar@{<=>}[r] \ar@{=>}[rd] \ar@{=>}[d] & (6) \\
& & (5) \ar@{=>}[u] & (4) \ar@{=>}[l] & (8)
}
$$

\medskip\noindent
Suppose that $M$ satisfies (1), (4), (5), (6), and (7).
We will prove that (3) holds. It suffices
to show that $M$ is projective. We have to show that $\Hom_R(M, -)$
is exact. Let $0 \to N'' \to N \to N'\to 0$ be a short exact sequence of
$R$-module. We have to show that
$0 \to \Hom_R(M, N'') \to \Hom_R(M, N) \to
\Hom_R(M, N') \to 0$ is exact.
As $M$ is finite locally free there exist a covering
$\Spec(R) = \bigcup D(f_i)$ such that $M_{f_i}$ is finite free.
By
Lemma \ref{lemma-hom-from-finitely-presented}
we see that
$$
0 \to \Hom_R(M, N'')_{f_i} \to \Hom_R(M, N)_{f_i} \to
\Hom_R(M, N')_{f_i} \to 0
$$
is equal to
$0 \to \Hom_{R_{f_i}}(M_{f_i}, N''_{f_i}) \to
\Hom_{R_{f_i}}(M_{f_i}, N_{f_i}) \to
\Hom_{R_{f_i}}(M_{f_i}, N'_{f_i}) \to 0$
which is exact as $M_{f_i}$ is free and as the localization
$0 \to N''_{f_i} \to N_{f_i} \to N'_{f_i} \to 0$
is exact (as localization is exact). Whence we see that
$0 \to \Hom_R(M, N'') \to \Hom_R(M, N) \to
\Hom_R(M, N') \to 0$ is exact by
Lemma \ref{lemma-cover}.

\medskip\noindent
Finally, assume that (8) holds. Pick a maximal ideal $\mathfrak m \subset R$.
Pick $x_1, \ldots, x_r \in M$ which map to a $\kappa(\mathfrak m)$-basis of
$M \otimes_R \kappa(\mathfrak m) = M/\mathfrak mM$. In particular
$\rho_M(\mathfrak m) = r$. By
Nakayama's Lemma \ref{lemma-NAK}
there exists an $f \in R$, $f \not \in \mathfrak m$ such that
$x_1, \ldots, x_r$ generate $M_f$ over $R_f$. By the assumption that
$\rho_M$ is locally constant there exists a $g \in R$, $g \not \in \mathfrak m$
such that $\rho_M$ is constant equal to $r$ on $D(g)$. We claim that
$$
\Psi : R_{fg}^{\oplus r} \longrightarrow M_{fg}, \quad
(a_1, \ldots, a_r) \longmapsto \sum a_i x_i
$$
is an isomorphism. This claim will show that $M$ is finite locally
free, i.e., that (7) holds. To see the claim
it suffices to show that the induced map on localizations
$\Psi_{\mathfrak p} : R_{\mathfrak p}^{\oplus r} \to M_{\mathfrak p}$
is an isomorphism for all $\mathfrak p \in D(fg)$, see
Lemma \ref{lemma-characterize-zero-local}.
By our choice of $f$ the map $\Psi_{\mathfrak p}$
is surjective. By assumption (8) we have
$M_{\mathfrak p} \cong R_{\mathfrak p}^{\oplus \rho_M(\mathfrak p)}$
and by our choice of $g$ we have $\rho_M(\mathfrak p) = r$.
Hence $\Psi_{\mathfrak p}$ determines a surjection
$R_{\mathfrak p}^{\oplus r} \to
M_{\mathfrak p} \cong R_{\mathfrak p}^{\oplus r}$
whence is an isomorphism by
Lemma \ref{lemma-fun}.
(Of course this last fact follows from a simple matrix argument also.)
\end{proof}

\begin{lemma}
\label{lemma-finite-projective-reduced}
Let $R$ be a reduced ring and let $M$ be an $R$-module. Then the
equivalent conditions of Lemma \ref{lemma-finite-projective}
are also equivalent to
\begin{enumerate}
\item[(9)] $M$ is finite and the function
$\rho_M : \Spec(R) \to \mathbf{Z}$, $\mathfrak p \mapsto
\dim_{\kappa(\mathfrak p)} M \otimes_R \kappa(\mathfrak p)$
is locally constant in the Zariski topology.
\end{enumerate}
\end{lemma}

\begin{proof}
Pick a maximal ideal $\mathfrak m \subset R$.
Pick $x_1, \ldots, x_r \in M$ which map to a $\kappa(\mathfrak m)$-basis of
$M \otimes_R \kappa(\mathfrak m) = M/\mathfrak mM$. In particular
$\rho_M(\mathfrak m) = r$. By
Nakayama's Lemma \ref{lemma-NAK}
there exists an $f \in R$, $f \not \in \mathfrak m$ such that
$x_1, \ldots, x_r$ generate $M_f$ over $R_f$. By the assumption that
$\rho_M$ is locally constant there exists a $g \in R$, $g \not \in \mathfrak m$
such that $\rho_M$ is constant equal to $r$ on $D(g)$. We claim that
$$
\Psi : R_{fg}^{\oplus r} \longrightarrow M_{fg}, \quad
(a_1, \ldots, a_r) \longmapsto \sum a_i x_i
$$
is an isomorphism. This claim will show that $M$ is finite locally
free, i.e., that (7) holds. Since $\Psi$ is surjective, it suffices to
show that $\Psi$ is injective. Since $R_{fg}$ is reduced,
it suffices to show that $\Psi$ is injective after localization
at all minimal primes $\mathfrak p$ of $R_{fg}$, see
Lemma \ref{lemma-reduced-ring-sub-product-fields}.
However, we know that $R_\mathfrak p = \kappa(\mathfrak p)$
by Lemma \ref{lemma-minimal-prime-reduced-ring} and
$\rho_M(\mathfrak p) = r$ hence $\Psi_\mathfrak p :
R_\mathfrak p^{\oplus r} \to M \otimes_R \kappa(\mathfrak p)$
is an isomorphism as a surjective map of finite dimensional
vector spaces of the same dimension.
\end{proof}

\begin{remark}
\label{remark-warning}
It is not true that a finite $R$-module which is
$R$-flat is automatically projective. A counter
example is where $R = \mathcal{C}^\infty(\mathbf{R})$
is the ring of infinitely differentiable functions on
$\mathbf{R}$, and $M = R_{\mathfrak m} = R/I$ where
$\mathfrak m = \{f \in R \mid f(0) = 0\}$ and
$I = \{f \in R \mid \exists \epsilon, \epsilon > 0 :
f(x) = 0\ \forall x, |x| < \epsilon\}$.
\end{remark}

\begin{lemma}
\label{lemma-finite-flat-local}
(Warning: see Remark \ref{remark-warning}.)
Suppose $R$ is a local ring, and $M$ is a finite
flat $R$-module. Then $M$ is finite free.
\end{lemma}

\begin{proof}
Follows from the equational criterion of flatness, see
Lemma \ref{lemma-flat-eq}. Namely, suppose that
$x_1, \ldots, x_r \in M$ map to a basis of
$M/\mathfrak mM$. By Nakayama's Lemma \ref{lemma-NAK}
these elements generate $M$. We want to show there
is no relation among the $x_i$. Instead, we will show
by induction on $n$ that if $x_1, \ldots, x_n \in M$
are linearly independent in the vector space
$M/\mathfrak mM$ then they are independent over $R$.

\medskip\noindent
The base case of the induction is where we have
$x \in M$, $x \not\in \mathfrak mM$ and a relation
$fx = 0$. By the equational criterion there
exist $y_j \in M$ and $a_j \in R$ such that
$x = \sum a_j y_j$ and $fa_j = 0$ for all $j$.
Since $x \not\in \mathfrak mM$ we see that
at least one $a_j$ is a unit and hence $f = 0 $.

\medskip\noindent
Suppose that $\sum f_i x_i$ is a relation among $x_1, \ldots, x_n$.
By our choice of $x_i$ we have $f_i \in \mathfrak m$.
According to the equational criterion of flatness there exist
$a_{ij} \in R$ and $y_j \in M$ such that
$x_i = \sum a_{ij} y_j$ and $\sum f_i a_{ij} = 0$.
Since $x_n \not \in \mathfrak mM$ we see that
$a_{nj}\not\in \mathfrak m$ for at least one $j$.
Since $\sum f_i a_{ij} = 0$ we get
$f_n = \sum_{i = 1}^{n-1} (-a_{ij}/a_{nj}) f_i$.
The relation $\sum f_i x_i = 0$ now can be rewritten
as $\sum_{i = 1}^{n-1} f_i( x_i + (-a_{ij}/a_{nj}) x_n) = 0$.
Note that the elements $x_i + (-a_{ij}/a_{nj}) x_n$ map
to $n-1$ linearly independent elements of $M/\mathfrak mM$.
By induction assumption we get that all the $f_i$, $i \leq n-1$
have to be zero, and also $f_n = \sum_{i = 1}^{n-1} (-a_{ij}/a_{nj}) f_i$.
This proves the induction step.
\end{proof}

\begin{lemma}
\label{lemma-finite-projective-descends}
Let $R \to S$ be a flat local homomorphism of local rings.
Let $M$ be a finite $R$-module. Then $M$ is finite projective
over $R$ if and only if $M \otimes_R S$ is finite projective
over $S$.
\end{lemma}

\begin{proof}
By Lemma \ref{lemma-finite-projective} being finite projective
over a local ring is the same thing as being finite free.
Suppose that $M \otimes_R S$ is a finite free $S$-module.
Pick $x_1, \ldots, x_r \in M$ whose images in $M/\mathfrak m_RM$
form a basis over $\kappa(\mathfrak m)$. Then
we see that $x_1 \otimes 1, \ldots, x_r \otimes 1$
are a basis for $M \otimes_R S$. This implies that
the map $R^{\oplus r} \to M, (a_i) \mapsto \sum a_i x_i$
becomes an isomorphism after tensoring with $S$.
By faithful flatness of $R \to S$, see Lemma \ref{lemma-local-flat-ff}
we see that it is an isomorphism.
\end{proof}

\begin{lemma}
\label{lemma-locally-free-semi-local-free}
Let $R$ be a semi-local ring.
Let $M$ be a finite locally free module.
If $M$ has constant rank, then $M$ is free. In particular, if $R$ has
connected spectrum, then $M$ is free.
\end{lemma}

\begin{proof}
Omitted. Hints: First show that $M/\mathfrak m_iM$ has the
same dimension $d$ for all maximal ideal $\mathfrak m_1, \ldots, \mathfrak m_n$
of $R$ using the rank is constant.
Next, show that there exist elements $x_1, \ldots, x_d \in M$
which form a basis for each $M/\mathfrak m_iM$ by the Chinese
remainder theorem. Finally show that $x_1, \ldots, x_d$ is a basis for $M$.
\end{proof}

\noindent
Here is a technical lemma that is used in the chapter on groupoids.

\begin{lemma}
\label{lemma-semi-local-module-basis-in-submodule}
Let $R$ be a local ring with maximal ideal $\mathfrak m$ and
infinite residue field.
Let $R \to S$ be a ring map.
Let $M$ be an $S$-module and let $N \subset M$ be an $R$-submodule.
Assume
\begin{enumerate}
\item $S$ is semi-local and $\mathfrak mS$ is contained in the
Jacobson radical of $S$,
\item $M$ is a finite free $S$-module, and
\item $N$ generates $M$ as an $S$-module.
\end{enumerate}
Then $N$ contains an $S$-basis of $M$.
\end{lemma}

\begin{proof}
Assume $M$ is free of rank $n$. Let $I \subset S$ be the Jacobson radical.
By Nakayama's Lemma \ref{lemma-NAK} a sequence of elements
$m_1, \ldots, m_n$ is a basis for $M$ if and only if
$\overline{m}_i \in M/IM$ generate $M/IM$. Hence we may replace
$M$ by $M/IM$, $N$ by $N/(N \cap IM)$, $R$ by $R/\mathfrak m$,
and $S$ by $S/IS$. In this case we see that $S$ is a finite product
of fields $S = k_1 \times \ldots \times k_r$ and
$M = k_1^{\oplus n} \times \ldots \times k_r^{\oplus n}$.
The fact that $N \subset M$ generates $M$ as an $S$-module
means that there exist $x_j \in N$ such that a linear combination
$\sum a_j x_j$ with $a_j \in S$ has a nonzero component in each
factor $k_i^{\oplus n}$.
Because $R = k$ is an infinite field, this means that also
some linear combination $y = \sum c_j x_j$ with $c_j \in k$ has a
nonzero component in each factor. Hence $y \in N$ generates a
free direct summand $Sy \subset M$. By induction on $n$ the result
holds for $M/Sy$ and the submodule $\overline{N} = N/(N \cap Sy)$.
In other words there exist $\overline{y}_2, \ldots, \overline{y}_n$
in $\overline{N}$ which (freely) generate $M/Sy$. Then
$y, y_2, \ldots, y_n$ (freely) generate $M$ and we win.
\end{proof}

\begin{lemma}
\label{lemma-evaluation-map-iso-finite-projective}
Let $R$ be ring. Let $L$, $M$, $N$ be $R$-modules.
The canonical map
$$
\Hom_R(M, N) \otimes_R L \to \Hom_R(M, N \otimes_R L)
$$
is an isomorphism if $M$ is finite projective.
\end{lemma}

\begin{proof}
By Lemma \ref{lemma-finite-projective} we see that $M$
is finitely presented as well as finite locally free.
By Lemmas \ref{lemma-hom-from-finitely-presented} and
\ref{lemma-tensor-product-localization} formation of
the left and right hand side of the arrow commutes with
localization. We may check that our map is an isomorphism
after localization, see Lemma \ref{lemma-cover}.
Thus we may assume $M$ is finite free. In this case
the lemma is immediate.
\end{proof}


 


